%% Question bank chapter file - generated by bank_tools.py build (do not edit while the bank is still being built:
%% it is regenerated from the per-test data in Question Bank\data). Every question is one bankq block:
%%   \begin{bankq}[id=..., source={...}, paper=.., qno=.., type=..., subtopic={...}, answer={...}, answer_src=key|solved,
%%                 working={...}, difficulty=1-10, diff_parts={K S I C F}, diff_note={...},
%%                 group=<passage id or empty>, also_in={duplicate copies in other tests}]
%% difficulty = min(10, K+S+I+C+F-1): K concepts 1-3, S steps 1-4, I insight 0-2, C computation 0-2, F format 0-1
%%   ...question LaTeX, options, figures (\bfig / \bfigi -> img/)...
%%   \end{bankq}
%% Passages shared by several questions are bankpassage blocks; their questions carry group=<passage id>.
\documentclass[10pt]{article}
\usepackage[a4paper,margin=1.8cm]{geometry}
\usepackage[T1]{fontenc}
\usepackage{amsmath,amssymb,mathtools}
\usepackage{graphicx}
\usepackage[export]{adjustbox}
\usepackage{xcolor,needspace,keyval}
\usepackage{enumitem}
\usepackage{booktabs,array,tabularx}
\usepackage[hidelinks]{hyperref}
\setlength{\parindent}{0pt}
\setlength{\parskip}{3pt}
\makeatletter
\define@key{bq}{id}{\def\bq@id{#1}}\define@key{bq}{source}{\def\bq@src{#1}}\define@key{bq}{paper}{\def\bq@paper{#1}}
\define@key{bq}{qno}{\def\bq@qno{#1}}\define@key{bq}{type}{\def\bq@type{#1}}\define@key{bq}{subtopic}{\def\bq@sub{#1}}
\define@key{bq}{answer}{\def\bq@ans{#1}}\define@key{bq}{answer_src}{\def\bq@asrc{#1}}\define@key{bq}{working}{\def\bq@work{#1}}
\define@key{bq}{group}{\def\bq@group{#1}}\define@key{bq}{also_in}{\def\bq@also{#1}}
\define@key{bq}{difficulty}{\def\bq@diff{#1}}\define@key{bq}{diff_parts}{\def\bq@dparts{#1}}\define@key{bq}{diff_note}{\def\bq@dnote{#1}}
\newenvironment{bankq}[1][]{\def\bq@work{}\def\bq@also{}\def\bq@group{}\def\bq@diff{?}\def\bq@dparts{}\def\bq@dnote{}\setkeys{bq}{#1}%
  \par\needspace{4\baselineskip}\medskip\noindent{\small\textbf{\bq@id}\hfill\color{gray}[\bq@src, Paper \bq@paper, Q\bq@qno]
  \bq@type\ $\cdot$ \bq@sub\ $\cdot$ \textbf{D\bq@diff}}\par\nopagebreak\relax}
  {\par{\small\color{teal}Answer: \textbf{\bq@ans}\ifx\bq@work\@empty\else\ (solved: \bq@work)\fi
  \ifx\bq@also\@empty\else\par\color{gray}Also in: \bq@also\fi}\par}
\define@key{bp}{id}{\def\bp@id{#1}}\define@key{bp}{source}{\def\bp@src{#1}}\define@key{bp}{questions}{\def\bp@qs{#1}}
\newenvironment{bankpassage}[1][]{\setkeys{bp}{#1}\par\needspace{4\baselineskip}\bigskip\noindent
  {\small\textbf{Passage \bp@id}\hfill\color{gray}[\bp@src; questions \bp@qs]}\par\nopagebreak\relax\itshape}{\par}
\makeatother
\newcommand{\bfig}[2]{\par\smallskip{\centering\includegraphics[width=#1pt,max width=0.62\linewidth,max height=0.32\textheight]{img/#2}\par}\smallskip}
\newcommand{\bfigi}[2]{\includegraphics[width=#1pt,max width=\linewidth,max height=0.32\textheight,valign=c]{img/#2}}
\newcommand{\bmissing}[1]{\textit{[figure missing in the source conversion: \detokenize{#1}]}}
\newenvironment{qoptions}{\begin{itemize}[leftmargin=1.8em,itemsep=1pt,topsep=2pt,label={}]}{\end{itemize}}
\newcommand{\opt}[2]{\item[(#1)] #2}
\providecommand{\fwstart}{}\providecommand{\fwend}{}\providecommand{\qfullsplit}{}
\newcolumntype{L}[1]{>{\raggedright\arraybackslash}p{#1}}
\newcolumntype{C}[1]{>{\centering\arraybackslash}p{#1}}
\newcolumntype{R}[1]{>{\raggedleft\arraybackslash}p{#1}}
\providecommand{\sectionhdr}[1]{}\providecommand{\subsectionhdr}[1]{}


\begin{document}
{\Large\bfseries Kinematics --- Question Bank}\par\smallskip
{\small Physics chapter P01. Questions from all converted test papers, grouped by source series. Answers: printed key where it matches the paper, otherwise solved (answer\_src=solved, with working). Exact repeats across papers appear once, with also\_in listing the other copies.}\par\medskip\hrule

\section*{Score-I papers}
\begin{bankpassage}[id=SI12-P2-PHY-S1-G7-8, source={Score-I - Test 12 (09 Mar 2022)}, questions={7-8}]
Column-I represent a particle of mass $m$ which projected as represent in column-I. In column-II a box of mass $m$ have its intial velocity $u$ \& acceleration $a$ as shown. In column-III a block of mass $m$ is put on wedge.
\begin{center}\small
\begin{tabular}{|p{0.8cm}|p{3.0cm}|p{0.8cm}|p{3.6cm}|p{0.8cm}|p{4.0cm}|}
\hline
 & \centering Column-I & & \centering Column-II & & \centering Column-III\tabularnewline
\hline
(I) & \bfigi{61.7}{SI12_p2_phy_s1_para7_8_1.png}\rule{0pt}{30pt} & (i) & \bfigi{70.8}{SI12_p2_phy_s1_para7_8_5.png} & (P) & \bfigi{56.8}{SI12_p2_phy_s1_para7_8_9.png}\\
\hline
(II) & \bfigi{63.9}{SI12_p2_phy_s1_para7_8_2.png}\rule{0pt}{30pt} & (ii) & \bfigi{81.7}{SI12_p2_phy_s1_para7_8_6.png} & (Q) & \bfigi{97.3}{SI12_p2_phy_s1_para7_8_10.png}\\
\hline
(III) & \bfigi{51.9}{SI12_p2_phy_s1_para7_8_3.png}\rule{0pt}{30pt} & (iii) & \bfigi{89.3}{SI12_p2_phy_s1_para7_8_7.png} & (R) & \bfigi{93.0}{SI12_p2_phy_s1_para7_8_11.png}\\
\hline
(IV) & \bfigi{57.0}{SI12_p2_phy_s1_para7_8_4.png}\rule{0pt}{30pt} & (iv) & \bfigi{70.5}{SI12_p2_phy_s1_para7_8_8.png} & (S) & \bfigi{93.5}{SI12_p2_phy_s1_para7_8_12.png}\\
\hline
\end{tabular}
\end{center}
\end{bankpassage}
\begin{bankq}[id=SI12-P2-PHY-S1-Q07, source={Score-I - Test 12 (09 Mar 2022)}, paper=2, qno=7, type=single, subtopic={Projectile motion}, answer={A}, answer_src=key, difficulty=2, diff_parts={K1 S2 I0 C0 F0}, diff_note={concepts: projectile motion cases | steps: vertical displacement for option A; option B; option C; option D}, group=SI12-P2-PHY-S1-G7-8]
In which combination displacement of particle in vertical direction is zero?
\begin{qoptions}
\opt{A}{(II) (i) (Q)}
\opt{B}{(III) (iii) (R)}
\opt{C}{(II) (iv) (S)}
\opt{D}{(II) (i) (P)}
\end{qoptions}
\end{bankq}
\begin{bankq}[id=SI12-P2-PHY-S1-Q08, source={Score-I - Test 12 (09 Mar 2022)}, paper=2, qno=8, type=single, subtopic={Projectile motion}, answer={B}, answer_src=key, difficulty=2, diff_parts={K1 S2 I0 C0 F0}, diff_note={concepts: projectile motion cases | steps: initial speed trend for option A; option B; option C; option D}, group=SI12-P2-PHY-S1-G7-8]
In which combination speed of particle must increase just after the motion start?
\begin{qoptions}
\opt{A}{(I) (i) (Q)}
\opt{B}{(III) (ii) (P)}
\opt{C}{(II) (iii) (R)}
\opt{D}{(I) (iv) (R)}
\end{qoptions}
\end{bankq}
\begin{bankq}[id=SI12-P2-PHY-S2-Q05, source={Score-I - Test 12 (09 Mar 2022)}, paper=2, qno=5, type=numerical, subtopic={Projectile motion}, answer={200.00}, answer_src=key, difficulty=4, diff_parts={K2 S1 I1 C0 F1}, diff_note={concepts: projectile motion, relative motion | steps: relative acceleration is zero so the relative path is a straight line; geometry of the collision point gives h | I: relative motion}]
Shots are fired from the top of a tower and from bottom simultaneously at angles $45^\circ$ and $60^\circ$ respectively as shown. If horizontal distance of the point of collision is at a distance of $(\sqrt3+1)$ m from the tower, then find the height $h$ of tower (in cm).
\bfig{107.2}{SI12_p2_phy_s2_q05_1.png}
\end{bankq}
\begin{bankq}[id=SI12-P2-PHY-S3-Q07, source={Score-I - Test 12 (09 Mar 2022)}, paper=2, qno=7, type=matching, subtopic={River-boat problems}, answer={P$\to$1, Q$\to$3, R$\to$2}, answer_src=key, difficulty=3, diff_parts={K1 S2 I0 C0 F1}, diff_note={concepts: relative velocity in river-boat motion | steps: resultant for 127 degrees; resultant for 90 degrees; resultant for 143 degrees}]
A river is flowing with speed 3 km/hr west to east. A man swims with speed 5 km/hr in still water. Man is at south bank of the river. Match the List-I with direction of velocities of man w.r.t. ground in List-II.
\begin{center}\footnotesize
\begin{tabular}{|p{0.6cm}|p{4.6cm}|p{0.6cm}|p{1.2cm}|}
\hline
 & \centering List-I & \multicolumn{2}{c|}{List-II}\\
\hline
(P) & Man swims at an angle $127^\circ$ from river flow. & (1) & \bfigi{6.3}{SI12_p2_phy_s3_q07_1.png}\rule{0pt}{16pt}\\
\hline
(Q) & Man swims right angle to river flow. & (2) & \bfigi{20.8}{SI12_p2_phy_s3_q07_2.png}\rule{0pt}{16pt}\\
\hline
(R) & Man swims at an angle $143^\circ$ from river flow. & (3) & \bfigi{21.5}{SI12_p2_phy_s3_q07_3.png}\rule{0pt}{16pt}\\
\hline
\end{tabular}
\end{center}
\end{bankq}

\section*{Score Advanced papers}
\begin{bankq}[id=SA01-P2-PHY-S2-Q10, source={Score Advanced - Test 01 (22 Sep 2020)}, paper=2, qno=10, type=single, subtopic={Relative motion}, answer={B}, answer_src=key, difficulty=7, diff_parts={K2 S2 I2 C2 F0}, diff_note={concepts: ants chasing each other, relative motion | steps: relative velocity; approach rate; time; path}]
Three ants are sitting on the vertices of an equilateral triangle of side of a. At t = 0, ant 1 starts approaching ant 2 with a speed v, ant 2 starts approaching ant 3 with a speed v and ant 3 start approaching ant 1 with a speed v. If a = 1m, v = 1m/s match the following\par
\noindent{\footnotesize\setlength{\tabcolsep}{2pt}\begin{tabular}{|L{.8cm}|L{4.2cm}|L{.8cm}|L{1.8cm}|}\hline
\multicolumn{2}{|c|}{List--I} & \multicolumn{2}{c|}{List--II}\\\hline
(P) & The three ant meet at the centroid of the triangle at time (in s) & (1) & $\dfrac{\sqrt3}2$\\\hline
(Q) & The angular velocity (in rad/s) of ant 2 with respect to ant 1 at t = 0 is & (2) & $\dfrac{25\sqrt3}{256}$\\\hline
(R) & The angular acceleration (in rad/s$^2$) of ant 2 with respect to ant 1 at t = 0 is & (3) & $\dfrac23$\\\hline
(S) & The area of the equilateral triangle (in m$^2$) formed by the three ants at any time t = a/4v will be & (4) & $\dfrac{3\sqrt3}4$\\\hline
\end{tabular}}\par
\begin{qoptions}
\opt{A}{P $\to$ 1 ; Q $\to$ 4 ; R $\to$ 2 ; S $\to$ 3}
\opt{B}{P $\to$ 3 ; Q $\to$ 1 ; R $\to$ 4 ; S $\to$ 2}
\opt{C}{P $\to$ 4 ; Q $\to$ 3 ; R $\to$ 1 ; S $\to$ 2}
\opt{D}{P $\to$ 3 ; Q $\to$ 1 ; R $\to$ 1 ; S $\to$ 4}
\end{qoptions}
\end{bankq}
\begin{bankq}[id=SA02-P1-PHY-S2-Q02, source={Score Advanced - Test 02 (13 May 2021)}, paper=1, qno=2, type=single, subtopic={1D and 2D motion}, answer={B}, answer_src=solved, working={Displacement graph rises then falls: speed decreases then increases with initial deceleration; x=t\^{}2+3t and the downward throw speed up continuously}, difficulty=3, diff_parts={K1 S2 I1 C0 F0}, diff_note={concepts: speed and acceleration of motions | steps: read the graph; sign of velocity and acceleration; match}]
In the list--I, there are certain situations depicted and list--II gives the description of speed and acceleration of particle. Match them. Neglect air resistance and assume uniform gravitational field wherever applicable. Consider time before it hits the ground.\par
\noindent{\footnotesize\setlength{\tabcolsep}{2pt}\begin{tabular}{|L{.7cm}|L{3.4cm}|L{.7cm}|L{2.7cm}|}\hline
 & List-I & & List--II\\\hline
(P) & \bfigi{103.0}{SA02_p1_phy_s2_q02_1.png}\par Particle moving in 1 dimension & (1) & speed increases \& then decreases\\\hline
(Q) & x = t$^2$ + 3t & (2) & speed decreases \& then increases\\\hline
(R) & A projectile thrown on level ground at an angle to the horizontal & (3) & speed increases continuously.\\\hline
(S) & A ball thrown vertically down from a cliff. & (4) & speed decreases continuously\\\hline
 & & (5) & acceleration is initially opposite to velocity.\\\hline
 & & (6) & acceleration is initially in the same direction as the velocity\\\hline
\end{tabular}}\par
\begin{qoptions}
\opt{A}{P $\to$ 1,5;Q $\to$ 2,6;R $\to$ 4;S $\to$ 2,5}
\opt{B}{P $\to$ 2,5;Q $\to$ 3,6;R $\to$ 2;S $\to$ 3,6}
\opt{C}{P $\to$ 1,5;Q $\to$ 4,6;R $\to$ 1,6;S $\to$ 2,5}
\opt{D}{P $\to$ 2,6;Q $\to$ 2;R $\to$ 3,6;S $\to$ 3,6}
\end{qoptions}
\end{bankq}
\begin{bankq}[id=SA03-P1-PHY-S1-Q01, source={Score Advanced - Test 03 (27 Jun 2021)}, paper=1, qno=1, type=multi, subtopic={Relative motion}, answer={A,C}, answer_src=solved, working={Ground wind is north (V); relative to the boat it is 53 deg N of W: V=16 m/s and relative speed 20 m/s}, difficulty=6, diff_parts={K2 S2 I1 C1 F1}, diff_note={concepts: wind seen from a boat and the shore, relative motion | steps: relative velocity triangle; wind speed from the ground flag; relative speed}]
A boat is traveling due east at 12 m/s. A flag on the boat flaps at 53$^\circ$ N of W. Another flag on the shore flaps due north.
\begin{qoptions}
\opt{A}{Speed of wind with respect to ground is 16 m/s}
\opt{B}{Speed of wind with respect to ground is 20 m/s}
\opt{C}{Speed of wind with respect to boat is 20 m/s}
\opt{D}{Speed of wind with respect to boat is 16 m/s}
\end{qoptions}
\end{bankq}
\begin{bankq}[id=SA03-P1-PHY-S3-Q04, source={Score Advanced - Test 03 (27 Jun 2021)}, paper=1, qno=4, type=matching, subtopic={1D and 2D motion}, answer={P$\to$3, Q$\to$1, R$\to$2}, answer_src=solved, working={t=4 s; separation 16 sqrt2 m; at t=8 relative speed 8 sqrt2}, difficulty=6, diff_parts={K2 S2 I1 C1 F1}, diff_note={concepts: two particles with perpendicular motions, 1d and 2d motion | steps: relative velocity zero time; positions; meeting time and relative speed}]
Two particles A and B moving in x-y plane are at origin at t = 0 sec. The initial velocity vectors of A and B are $\vec{u}_A = 8\hat{i}$ m/s and $\vec{u}_B = 8\hat{j}$ m/s. The acceleration of A and B are constant and are $\vec{a}_A = -2\hat{i}$ m/s$^2$ and $\vec{a}_B = -2\hat{j}$ m/s$^2$. List-I gives certain statements regarding particle A and B. List-II gives corresponding results. Match the statements in List-I with corresponding results in List-II.\par
\noindent{\footnotesize\setlength{\tabcolsep}{2pt}\begin{tabular}{|L{.7cm}|L{4.0cm}|L{.7cm}|L{1.3cm}|}\hline
 & List-I & & List-II\\\hline
(P) & The time (in seconds) at which velocity of A relative to B is zero & (1) & 16$\sqrt{2}$\\\hline
(Q) & The distance (in metres) between A and B when their relative velocity is zero & (2) & 8$\sqrt{2}$\\\hline
(R) & The magnitude of relative velocity of A and B at the instant they are at same position. & (3) & 4\\\hline
\end{tabular}}\par
\end{bankq}
\begin{bankq}[id=SA03-P2-PHY-S1-Q09, source={Score Advanced - Test 03 (27 Jun 2021)}, paper=2, qno=9, type=single, subtopic={Projectile motion}, answer={D}, answer_src=solved, working={dv\_x=0 and dv\_y=-g dt}, difficulty=1, diff_parts={K1 S1 I0 C0 F0}, diff_note={concepts: velocity changes of a projectile | steps: horizontal constant; vertical decreases}]
With the x-axis horizontal and the y-axis vertically upward as positive, the change in the horizontal component of velocity, $\Delta v_x$, and the change in the vertical component of velocity, $\Delta v_y$, of a projectile are related to the time since leaving the barrel, $\Delta t$, as
\begin{qoptions}
\opt{A}{$\Delta v_x = 0;\ \Delta v_y = 0$}
\opt{B}{$\Delta v_x = g\Delta t;\ \Delta v_y = 0$}
\opt{C}{$\Delta v_x = 0;\ \Delta v_y = g\Delta t$}
\opt{D}{$\Delta v_x = 0;\ \Delta v_y = -g\Delta t$}
\end{qoptions}
\end{bankq}
\begin{bankq}[id=SA03-P2-PHY-S2-Q02, source={Score Advanced - Test 03 (27 Jun 2021)}, paper=2, qno=2, type=numerical, subtopic={Relative motion}, answer={6.25}, answer_src=solved, working={u=5, w=3.75 so 6.25 km/h}, difficulty=7, diff_parts={K2 S2 I2 C1 F1}, diff_note={concepts: rain speed from two apparent angles, relative motion | steps: components of rain; condition for equal angle; speed}]
Rain appears to be falling at 53$^\circ$ with the vertical to a man standing on the ground. When he starts running at 10 km/hr, rain again appears to be falling at 53$^\circ$ with the vertical. Find the speed (in km/hr) of rain wrt ground :-
\end{bankq}
\begin{bankq}[id=SA03-P2-PHY-S2-Q03, source={Score Advanced - Test 03 (27 Jun 2021)}, paper=2, qno=3, type=numerical, subtopic={Projectile motion}, answer={75}, answer_src=solved, working={T=5 s, range=(1/2)(6)(25)=75 m}, difficulty=7, diff_parts={K2 S2 I2 C1 F1}, diff_note={concepts: range on an incline for a perpendicular throw, projectile motion | steps: axes along and perpendicular to the incline; time of flight; range}]
A projectile is thrown perpendicular to a long inclined plane as shown. The range (in m) of projectile on the incline is :-\par
\bfig{79.0}{SA03_p2_phy_s2_q03_1.png}
\end{bankq}
\begin{bankq}[id=SA03-P2-PHY-S2-Q07, source={Score Advanced - Test 03 (27 Jun 2021)}, paper=2, qno=7, type=numerical, subtopic={Projectile motion}, answer={10}, answer_src=solved, working={1875+25(25-10t)=0 gives t=10 s}, difficulty=3, diff_parts={K1 S1 I1 C0 F1}, diff_note={concepts: time when velocity is perpendicular to the initial one | steps: dot product zero}]
A ball is thrown with 50 m/s at 30$^\circ$ with the horizontal from a building of height 500 m. The time (in sec) at which velocity vector of the ball becomes perpendicular to the initial velocity is (g = 10 m/s$^2$) :-
\end{bankq}
\begin{bankq}[id=SA03-P2-PHY-S2-Q10, source={Score Advanced - Test 03 (27 Jun 2021)}, paper=2, qno=10, type=numerical, subtopic={Projectile motion}, answer={312.5}, answer_src=solved, working={v=50, a\_n=8, R=312.5 m}, difficulty=5, diff_parts={K1 S2 I1 C1 F1}, diff_note={concepts: radius of curvature of a trajectory | steps: velocity at 7 s; normal acceleration; radius}]
A ball is thrown with $40\sqrt{2}$ m/s at 45$^\circ$ with horizontal. Find the radius of curvature (in m) of its path at t = 7 sec :-
\end{bankq}
\begin{bankq}[id=SA04-P2-PHY-S1-Q12, source={Score Advanced - Test 04 (04 Jul 2021)}, paper=2, qno=12, type=single, subtopic={Projectile motion}, answer={B}, answer_src=solved, working={Straight line, time of flight 2 s, displacement zero. The range (horizontal displacement) is 0, not 20 m}, difficulty=6, diff_parts={K2 S2 I2 C1 F0}, diff_note={concepts: charge projected in horizontal field and gravity, projectile motion | steps: acceleration (-10,-10) is antiparallel to the velocity (10,10); straight line motion; time of flight}]
A particle of mass 1gm and charge $-$0.1$\mu$C is projected from ground with a velocity 10$\sqrt{2}$ at an 45$^\circ$ with horizontal in the area having uniform electric field 1 kV/cm in horizontal direction. Acceleration due to gravity is 10m/s$^2$ in vertical downward direction. Select incorrect statement:\par
\bfig{131.2}{SA04_p2_phy_s1_q12_1.png}
\begin{qoptions}
\opt{A}{Time of flight of particle is 2 sec}
\opt{B}{Range of particle is 20 m}
\opt{C}{Total displacement of particle is 0m}
\opt{D}{Particle will follow straight line motion.}
\end{qoptions}
\end{bankq}
\begin{bankq}[id=SA05-P1-PHY-S2-Q04, source={Score Advanced - Test 05 (15 Jul 2021)}, paper=1, qno=4, type=numerical, subtopic={1D and 2D motion}, answer={0.25}, answer_src=key, difficulty=5, diff_parts={K2 S1 I1 C1 F1}, diff_note={concepts: speed at a time for a curved path, 1d and 2d motion | steps: x=t\^{}2/2, y=t\^{}4/8; differentiate}]
A particle moves along a curve y = $\dfrac{x^2}{2}$ and x = $\dfrac{t^2}{2}$, where x and y are measured in meter and t is in seconds. Then the speed of particle at t = 2 sec is v m/s then $\dfrac{v^2}{80}$ is
\end{bankq}
\begin{bankq}[id=SA06-P1-PHY-S3-Q07, source={Score Advanced - Test 06 (04 Sep 2021)}, paper=1, qno=7, type=numerical, subtopic={River-boat problems}, answer={4}, answer_src=key, difficulty=8, diff_parts={K3 S1 I2 C2 F1}, diff_note={concepts: boat crossing along a fixed line and back, river-boat problems, resolve velocities relative to the river for each direction | steps: resolve velocities relative to the river for each direction; sum of times}]
A boat travel, between two points A and B on the opposite banks of a river (figure), always following the line AB. The distance S between points A and B is 240 m. The velocity of the river flow v$_r$ = $2\sqrt{2}$ m/s is constant over the entire width of the river. The line AB makes an $\alpha$ = 45$^\circ$ with the direction of the current. The boat moves to cover the distance AB and back in a time t = $2\sqrt{3}$ min. The angle $\beta$ remains the same during the passage from A to B and from B to A. With what velocity u$_b$ should the boat travel ?\par
\bfig{122.0}{SA06_p1_phy_s3_q07_1.png}
\end{bankq}
\begin{bankq}[id=SA07-P1-PHY-S1-Q06, source={Score Advanced - Test 07 (16 Apr 2023)}, paper=1, qno=6, type=multi, subtopic={1D and 2D motion}, answer={A,C}, answer_src=key, difficulty=6, diff_parts={K2 S1 I2 C1 F1}, diff_note={concepts: speeding up or slowing down when acceleration points towards a point, 1d and 2d motion | steps: tangential component of acceleration along the velocity}]
A small object moves counter clockwise along the circular path whose centre is at origin as shown in figure. As it moves along the path, its acceleration vector continuously points towards point S. Then the object\par
\bfig{104.4}{SA07_p1_phy_s1_q06_1.png}
\begin{qoptions}
\opt{A}{Speed up as it moves from A to C via B.}
\opt{B}{Slows down as it moves from A to C via B.}
\opt{C}{Slows down as it moves from C to A via D.}
\opt{D}{Speed up as it moves from C to A via D.}
\end{qoptions}
\end{bankq}
\begin{bankq}[id=SA08-P2-PHY-S3-Q06, source={Score Advanced - Test 08 (22 Apr 2023)}, paper=2, qno=6, type=numerical, subtopic={1D and 2D motion}, answer={1}, answer_src=key, difficulty=8, diff_parts={K3 S1 I2 C2 F1}, diff_note={concepts: pursuit along a radius, 1d and 2d motion, radial and angular velocity components | steps: radial and angular velocity components; time of chase}]
A boy runs on a circular path of radius R = 28 m with a constant speed u = 4 m/s. Another boy starts from the centre of the path to catch the first boy. The second boy always remains on the radius connecting the centre of the circle and the first boy and maintains magnitude of his velocity constant V = 4 m/s. If the time of chase is (10 + x) sec then x =.....; $\left(\pi = \dfrac{22}{7}\right)$
\end{bankq}
\begin{bankq}[id=SA10-P1-PHY-S3-Q05, source={Score Advanced - Test 10 (06 May 2023)}, paper=1, qno=5, type=numerical, subtopic={Projectile motion}, answer={240}, answer_src=key, difficulty=8, diff_parts={K3 S1 I2 C2 F1}, diff_note={concepts: difference of ranges of balls thrown up and down from a cliff, projectile motion, time of flight for each | steps: time of flight for each; horizontal range difference}]
Two balls are thrown with the same speed 50 m/s from the top of a cliff. The angles of their initial velocitites are 37\textdegree{} above and below the horizontal, as shown in figure. How much farther (in m) along the ground does the top ball hit than the bottom ball?\par
\bfig{92.6}{SA10_p1_phy_s3_q05_1.png}
\end{bankq}
\begin{bankq}[id=SA10-P2-PHY-S2-Q07, source={Score Advanced - Test 10 (06 May 2023)}, paper=2, qno=7, type=multi, subtopic={1D and 2D motion}, answer={B,D}, answer_src=key, difficulty=8, diff_parts={K3 S1 I2 C2 F1}, diff_note={concepts: motion from velocity time graphs, 1d and 2d motion, area under the graphs and slopes | steps: area under the graphs and slopes}]
Two particles A and B start to move from origin along +ve x-axis. In the time interval t = 0 to t = 2 sec. the velocity -time graph of both particles is shown in below figures.\par
\bfig{260.3}{SA10_p2_phy_s2_q07_1.png}

\begin{qoptions}
\opt{A}{Separation between A and B is maximum at t = 1 sec.}
\opt{B}{Separation between A and B is maximum at t = 1.5 sec.}
\opt{C}{Maximum Separation between A and B is 2m}
\opt{D}{Maximum Separation between A and B is 2.5m}
\end{qoptions}
\end{bankq}
\begin{bankq}[id=SA11-P1-PHY-S1-Q01, source={Score Advanced - Test 11 (09 May 2023)}, paper=1, qno=1, type=multi, subtopic={Projectile motion}, answer={B,C,D}, answer_src=key, difficulty=8, diff_parts={K3 S1 I2 C2 F1}, diff_note={concepts: two projectiles launched together that collide, projectile motion, equal horizontal and vertical positions at the collision | steps: equal horizontal and vertical positions at the collision}]
Suppose two particles 1 and 2 are projected in vertical plane simultaneously.\par
\bfig{184.4}{SA11_p1_phy_s1_q01_1.png}
Their angles of projection are 30\textdegree{} and $\theta$ , respectively, with the horizontal. Let they collide after a time t in air. Then (neglect air resistance)
\begin{qoptions}
\opt{A}{$\theta$ = sin$^{-1}$(4/5) and they will have same speed just before the collision.}
\opt{B}{$\theta$ = sin$^{-1}$(4/5) and they will have different speed just before the collision.}
\opt{C}{x $<$1280$\sqrt{3}-$960m}
\opt{D}{It is possible that the particles collide when both of them are at their highest point.}
\end{qoptions}
\end{bankq}
\begin{bankq}[id=SA12-P1-PHY-S1-Q01, source={Score Advanced - Test 12 (11 May 2023)}, paper=1, qno=1, type=single, subtopic={Relative motion}, answer={A}, answer_src=key, difficulty=5, diff_parts={K2 S1 I2 C1 F0}, diff_note={concepts: collinearity condition for three moving points, relative motion | steps: y of the middle point is the average of the other two; differentiate twice}, also_in={Nurture (General) - Test 01 (12 Jun 2021) P2 Q2}]
At the initial moment three point A, B and C are on a horizontal plane along a straight line such that AB = BC. Point A begins to move upward with a constant velocity `v' and point C downward without any initial velocity at a constant acceleration `a'. If the points being to move simultaneously, then the initial velocity and acceleration of point B for all the three particles to be constantly on same straight line must be
\begin{qoptions}
\opt{A}{$\dfrac{v}{2}$ upwards, $\dfrac{a}{2}$ downwards}
\opt{B}{$\dfrac{v}{2}$ upwards, $\dfrac{a}{2}$ upwards}
\opt{C}{$\dfrac{v}{2}$ downwards, $\dfrac{a}{2}$ downwards}
\opt{D}{$\dfrac{v}{2}$ downwards, $\dfrac{a}{2}$ upwards}
\end{qoptions}
\end{bankq}
\begin{bankq}[id=SA12-P1-PHY-S3-Q06, source={Score Advanced - Test 12 (11 May 2023)}, paper=1, qno=6, type=numerical, subtopic={River-boat problems}, answer={9}, answer_src=key, difficulty=8, diff_parts={K3 S1 I2 C2 F1}, diff_note={concepts: crossing a channel using a plank, river-boat problems, relative motion and the time of crossing | steps: relative motion and the time of crossing}]
A water channel of width b = 17 m lies between two frictionless platforms A and B. A plank (P) of length $\ell$ = 2 m rest near the platform A. A block of mass m = 25 kg sliding on the platform A lands on the plank and stays on it. Force $\vec{\text{F}}$ of water resistance on the plank is proportional to velocity of the plank ($\vec{\text{v}}_\text{pr}$) relative to the water and given by the law $\vec{\text{F}}$ = $-$k$\vec{\text{v}}_\text{pr}$, where k = 15N -s/m. Calculate the minimum speed u$_\text{min}$ of the block on platform A so that plank reaches the platform B. (Answer in m /s)\par
\bfig{135.1}{SA12_p1_phy_s3_q06_1.png}
\end{bankq}
\begin{bankq}[id=SA13-P1-PHY-S1-Q01, source={Score Advanced - Test 13 (13 May 2023)}, paper=1, qno=1, type=numerical, subtopic={1D and 2D motion}, answer={60}, answer_src=key, difficulty=8, diff_parts={K3 S1 I2 C2 F1}, diff_note={concepts: motion of a supersonic jet relative to the ground, 1d and 2d motion, relative speed and heading of the plane | steps: relative speed and heading of the plane; distance in a given time}]
In 1997, Indian Airforce's most secret supersonic plane `Foxbat Jet', which can fly upto a maximum height of 23000 m and have maximum velocity of 3500 kmph, was flying in Pakistani Airspace at a height of 20,000 m with subsonic velocity. At some point of time, pilot of the plane decided to go supersonic and will fly with a speed of 1020 m/s, this created a sonic boom. Assuming velocity of sound to be 340 m/s. At what distance (in km) from border pilot should have attempted this to safely reach Indian Airspace? Assume he gets noticed when sound reaches the ground.\par
\end{bankq}
\begin{bankpassage}[id=SA16-P1-PHY-S3-G5-6, source={Score Advanced - Test 16 (20 May 2023)}, questions={5-6}]
\textbf{Question Stem}\par
A particle is projected from a point P (0, 0, 600 m) lying in horizontal X - Z plane with initial velocity of $\vec u=\left(80\hat i+100\hat j-60\hat k\right)$ m/s at $t=0$. It hits a smooth vertical wall in X - Y plane with y-axis in vertically upward direction. Co-efficient of restitution $e=\dfrac12$. (g = 10 m/s$^2$)\par
\end{bankpassage}
\begin{bankq}[id=SA16-P1-PHY-S3-Q05, source={Score Advanced - Test 16 (20 May 2023)}, paper=1, qno=5, type=numerical, subtopic={Projectile motion}, answer={9}, answer_src=key, difficulty=9, diff_parts={K3 S2 I2 C2 F1}, diff_note={concepts: angle with vertical when a projectile hits a wall, projectile motion, time to reach the wall from x | steps: time to reach the wall from x; velocity components; angle}, group=SA16-P1-PHY-S3-G5-6]
The ball hits the wall at time $t$ (in sec) when its velocity is making an angle $\theta$ (in degrees) with the vertical. Then write $\theta/t$
\end{bankq}
\begin{bankq}[id=SA16-P1-PHY-S3-Q06, source={Score Advanced - Test 16 (20 May 2023)}, paper=1, qno=6, type=numerical, subtopic={Projectile motion}, answer={1900}, answer_src=key, difficulty=9, diff_parts={K3 S2 I2 C2 F1}, diff_note={concepts: landing point after a restitution collision with a wall, collisions, x velocity reverses with e | steps: x velocity reverses with e; time of flight; coordinates}, group=SA16-P1-PHY-S3-G5-6]
Position of the point of impact with the horizontal X-Z plane after collision with the wall is (x,y,z). Find (x + y + z).
\end{bankq}
\begin{bankq}[id=SA17-P1-PHY-S3-Q06, source={Score Advanced - Test 17 (25 May 2023)}, paper=1, qno=6, type=numerical, subtopic={Relative motion}, answer={4.50}, answer_src=key, difficulty=6, diff_parts={K2 S1 I2 C1 F1}, diff_note={concepts: velocity of the middle of a pulled tape, relative motion | steps: the moving half has a velocity changing along the length}]
We stick a inextensible adhesive tape on a tabletop. We grab one end of it and pull it back horizontally and uniformly with a velocity of 6 cm/s. What is the velocity (in cm/s) of the middle of the moving part of the tape?\par
\bfig{103.1}{SA17_p1_phy_s3_q06_1.png}
\end{bankq}

\section*{Enthusiast Pre-Score papers}
\begin{bankq}[id=EPS01-P1-PHY-S2-Q07, source={Enthusiast Pre-Score - Test 01 (04 Nov 2020)}, paper=1, qno=7, type=multi, subtopic={Projectile motion}, answer={A,D}, answer_src=key, difficulty=9, diff_parts={K2 S3 I2 C2 F1}, diff_note={concepts: projectile with wind acceleration, projectile motion | steps: horizontal and vertical motions; times of equal height; times of equal x; ON; OM}]
A particle is projected from origin with speed u=50 m/s at an angle 30$^\circ$ with X-axis. Due to wind particle acquires an extra acceleration a = 6 m/s$^2$ in $-$ve x-direction. Path of particle is shown in figure.\par
\bfig{237.8}{EPS01_p1_phy_s2_q07_1.png}
x - Coordinates of point A and B are same\par
y- Co-ordinates of P and Q are same.\par
Point N on the X-axis, point M on the Y-axis.\par
Time taken by particle from O to P, O to Q, O to A and O to B are t$_1$, t$_2$, t$_3$ and t$_4$ respectively. (Assuming Y along vertical and their is no surface in the space) (g = 10 m/s$^2$ and constant).
\begin{qoptions}
\opt{A}{The value of t$_1$ + t$_2$ is 5 sec.}
\opt{B}{The value of t$_3$ + t$_4$ is 25 sec.}
\opt{C}{The length of ON is $125\sqrt3\,(5-\sqrt3)$ m}
\opt{D}{The length of OM is $\dfrac{625}{\sqrt3}\left(\dfrac5{\sqrt3}-1\right)$ m}
\end{qoptions}
\end{bankq}
\begin{bankq}[id=EPS01-P1-PHY-S2-Q09, source={Enthusiast Pre-Score - Test 01 (04 Nov 2020)}, paper=1, qno=9, type=multi, subtopic={Projectile motion}, answer={A,D}, answer_src=key, difficulty=8, diff_parts={K2 S2 I2 C2 F1}, diff_note={concepts: projectile meeting a sliding block, projectile motion | steps: components along the incline; time of meeting; height; collision}]
A particle of mass m is projected up from the bottom of an smooth inclined plane with initial velocity v$_0$ at angle 45$^\circ$ with an inclined plane of inclination 30$^\circ$ as shown in figure. At the same time, a small block of same mass m is released from rest at a height h. The particle hits the block at some point on inclined plane. (take g = 10 m/s$^2$) (size of the block is negligible)\par
\bfig{126.5}{EPS01_p1_phy_s2_q09_1.png}
\begin{qoptions}
\opt{A}{The value of height h is $\dfrac{v_0^2}{\sqrt3\mathrm g}$\,.}
\opt{B}{The value of height h is $\dfrac{\sqrt3v_0^2}{\mathrm g}$\,.}
\opt{C}{If the particle stick to the block after collision, the velocity of block parallel to the inclined plane just after collision will be $\dfrac{v_0}{\sqrt2}$\,.}
\opt{D}{If the particle stick to the block after collision, the velocity of block parallel to the inclined plane just after collision will be $\dfrac{v_0}{2\sqrt2}\left(\dfrac4{\sqrt3}-1\right)$\,.}
\end{qoptions}
\end{bankq}
\begin{bankq}[id=EPS01-P2-PHY-S1-Q05, source={Enthusiast Pre-Score - Test 01 (04 Nov 2020)}, paper=2, qno=5, type=multi, subtopic={Relative motion}, answer={A,B,D}, answer_src=key, difficulty=8, diff_parts={K2 S2 I2 C2 F1}, diff_note={concepts: particle meeting a moving wedge normally, relative motion | steps: relative motion; normal strike condition; time and distances; path shape}]
A large wedge BCD having its inclined surface at an angle $\theta$ = 45$^\circ$ to the horizontal, is travelling horizontally leftwards with uniform velocity v$_1$ = 10 m/s. At some instant, a particle is projected vertically up with speed v$_2$ = 20 m/s from point A on ground lying at some distance left to the lower edge B of the wedge as shown in figure. The particle strikes the incline BC normally, while it was falling.\par
\bfig{188.5}{EPS01_p2_phy_s1_q05_1.png}
\begin{qoptions}
\opt{A}{The distance AB at the instant the particle was projected from A, is 15 m.}
\opt{B}{The distance of lower edge B of the wedge from point A at the instant particle strikes the incline, is 15 m.}
\opt{C}{The distance of lower edge B of the wedge from point A at the instant particle strikes the incline, is 30 m.}
\opt{D}{The path of the particle in the reference frame attached to wedge is parabolic.}
\end{qoptions}
\end{bankq}

\section*{TOAS papers}
\begin{bankq}[id=TOAS02-P1-PHY-S1-Q01, source={TOAS - Test 02 (13 Feb 2022)}, paper=1, qno=1, type=single, subtopic={1D and 2D motion}, answer={B}, answer_src=key, difficulty=2, diff_parts={K1 S2 I0 C0 F0}, diff_note={concepts: motion with a velocity proportional to depth | steps: y = g t\^{}2/2 gives the fall time 4 s; x = integral of a y dt with y(t); evaluate at t = 4 s}]
A particle is released from rest from a certain height $H=80$ m. Due to the wind the particle gathers the horizontal velocity component $v_x=ay$ where $a=3/8\ \mathrm{s^{-1}}$ and $y$ is the vertical displacement of the particle from the point of release. Find the horizontal displacement of the particle when it strikes the ground. (Take : $g=10\ \mathrm{m/s^2}$)
\begin{qoptions}
\opt{A}{50 m}
\opt{B}{40 m}
\opt{C}{30 m}
\opt{D}{None of these}
\end{qoptions}
\end{bankq}
\begin{bankq}[id=TOAS02-P1-PHY-S1-Q06, source={TOAS - Test 02 (13 Feb 2022)}, paper=1, qno=6, type=single, subtopic={Projectile motion}, answer={B}, answer_src=key, difficulty=2, diff_parts={K1 S1 I1 C0 F0}, diff_note={concepts: region reachable by a projectile | steps: envelope of trajectories (safety parabola) for u\^{}2 = 2 g alpha; target outside means 4 beta > 3 alpha | I: safety parabola}]
A projectile is required to hit a target whose co-ordinates relative to horizontal and vertical axes through the point of projection are $(\alpha,\beta)$. If the gun velocity is $\sqrt{2g\alpha}$, then the condition when it is impossible to hit the target, is
\begin{qoptions}
\opt{A}{$3\alpha>4\beta$}
\opt{B}{$4\beta>3\alpha$}
\opt{C}{$3\beta>2\alpha$}
\opt{D}{$4\alpha>3\beta$}
\end{qoptions}
\end{bankq}
\begin{bankq}[id=TOAS02-P1-PHY-S2-Q10, source={TOAS - Test 02 (13 Feb 2022)}, paper=1, qno=10, type=multi, subtopic={Relative motion}, answer={A,D}, answer_src=key, difficulty=3, diff_parts={K1 S2 I0 C0 F1}, diff_note={concepts: relative motion | steps: positions of both boats at a common time; relative velocity; minimum separation; time of minimum separation}]
Two boats are at sea and moving continuously in straight lines. At 11 AM boat 1 was 40 km North of a small island and was moving at 15 km/h towards east. The boat 2 at 7 AM of the same day was 100 km away east of the same island and moved at a speed of 15 km/h towards the South. Which of the following is (are) \textbf{CORRECT}?
\begin{qoptions}
\opt{A}{The minimum separation between the boats is $100\sqrt2$ km}
\opt{B}{The minimum separation between the boats is less than $100\sqrt2$ km}
\opt{C}{The boats will be found at minimum separation at 9 AM}
\opt{D}{The boats will be found at minimum separation at 11 AM}
\end{qoptions}
\end{bankq}
\begin{bankq}[id=TOAS02-P1-PHY-S2-Q11, source={TOAS - Test 02 (13 Feb 2022)}, paper=1, qno=11, type=multi, subtopic={Projectile motion}, answer={A,C,D}, answer_src=key, difficulty=3, diff_parts={K1 S2 I0 C0 F1}, diff_note={concepts: projectile motion on a ramp | steps: components of velocity; trajectory equation; landing on the ramp or plateau in options A and B; displacement in options C and D}]
In the figure shown, a ball is launched with a velocity of magnitude 10 m/s at an angle $53^\circ$ with the horizontal. The launch point is at the base of a ramp of horizontal length $d_1$ and height $d_2$. A plateau (flat surface) is located at the top of the ramp. Which of the following option(s) is(are) \textbf{CORRECT}? (Take : $g=10\ \mathrm{m/s^2}$)
\bfig{207.5}{TOAS02_p1_phy_s2_q11_1.png}
\begin{qoptions}
\opt{A}{If $d_1=12$ m and $d_2=5$m, the ball lands of ramp.}
\opt{B}{If $d_1=6$m and $d_2=5$m, ball lands on plateau.}
\opt{C}{If $d_1=6$m and $d_2=3$m, the displacement of the ball from the point of projection to the point where it hits the surface is $3\sqrt5$m}
\opt{D}{If $d_1=6$m and $d_2=1$m, displacement of the particle is $>3\sqrt5$m}
\end{qoptions}
\end{bankq}
\begin{bankq}[id=TOAS02-P1-PHY-S3-Q02, source={TOAS - Test 02 (13 Feb 2022)}, paper=1, qno=2, type=numerical, subtopic={1D and 2D motion}, answer={162}, answer_src=key, difficulty=6, diff_parts={K1 S3 I1 C1 F1}, diff_note={concepts: relative motion of a bird and a car | steps: first trip of the bird to A and the car position; meeting time with the car; return to A; repeat for the second cycle; find the ratio of successive gaps | I: geometric pattern of the trips}]
A bird flies with a uniform speed of 10 km/h and a car moves with a uniform speed of 8 km /h. Both start from $B$ towards $A$ ($BA=40$ km) at the same instant. The bird having reached $A$, flies back to meet the approaching car. As soon as it reaches the car, it flies back to $A$. Assume that bird immediately gains the same speed while turning back. The distance of the car from point $A$ is $\dfrac{16}{\alpha}$ when the fly reaches point $A$ for the 3rd time. The value of $\alpha$ is :
\end{bankq}
\begin{bankq}[id=TOAS03-P1-PHY-S1-Q06, source={TOAS - Test 03 (27 Feb 2022)}, paper=1, qno=6, type=single, subtopic={1D and 2D motion}, answer={D}, answer_src=key, difficulty=2, diff_parts={K1 S1 I1 C0 F0}, diff_note={concepts: acceleration from a 1/V versus t graph | steps: a = -V\^{}2 d(1/V)/dt; read V and the slope at t = 3 s | I: chain rule on 1/V}]
The graph shows the variation of $\dfrac1V$ (where $V$ is the velocity of the particle) with respect to time. Then find the value of acceleration at $t=3$ sec.
\bfig{124.9}{TOAS03_p1_phy_s1_q06_1.png}
\begin{qoptions}
\opt{A}{$5\ \mathrm{m/s^2}$}
\opt{B}{$1\ \mathrm{m/s^2}$}
\opt{C}{$3\sqrt3\ \mathrm{m/s^2}$}
\opt{D}{$3\ \mathrm{m/s^2}$}
\end{qoptions}
\end{bankq}
\begin{bankq}[id=TOAS03-P1-PHY-S2-Q12, source={TOAS - Test 03 (27 Feb 2022)}, paper=1, qno=12, type=multi, subtopic={1D and 2D motion}, answer={C,D}, answer_src=key, difficulty=2, diff_parts={K1 S1 I0 C0 F1}, diff_note={concepts: motion with velocity proportional to displacement | steps: a = v dv/dx = 8v is positive; integrate dx/(8x + 1) = dt to get x(t)}]
The instantaneous velocity $v$ of a particle is related to its displacement $x$ according to the relation $v=8x+1$ (All parameters in SI units). Which of the following statement(s) is (are) true if $x=0$ at $t=0$ :-
\begin{qoptions}
\opt{A}{The particle will experience a retardation.}
\opt{B}{The particle will be at rest at $t=0$.}
\opt{C}{The displacement of the particle at time $t$ is $x=\dfrac18\left(e^{8t}-1\right)$}
\opt{D}{The particle will experience an acceleration.}
\end{qoptions}
\end{bankq}
\begin{bankq}[id=TOAS03-P1-PHY-S3-Q06, source={TOAS - Test 03 (27 Feb 2022)}, paper=1, qno=6, type=numerical, subtopic={Relative motion}, answer={1.50}, answer_src=key, difficulty=2, diff_parts={K1 S1 I0 C0 F1}, diff_note={concepts: closest approach by relative motion | steps: relative position and velocity; t = -(r.v)/v\^{}2}]
Two particles $A$ and $B$ located at points $(0,-10\sqrt3)$ and $(0,0)$ respectively in $xy$ plane. They start moving simultaneously at time $t=0$ with constant velocities $\vec V_A=5\hat i$ m/s and $\vec V_B=-5\sqrt3\,\hat j$ m/s, respectively. Time when they are closest to each other is found to be $K$ second. Find $K$. All distance are given in meter.
\end{bankq}
\begin{bankq}[id=TOAS03-P2-PHY-S2-Q01, source={TOAS - Test 03 (27 Feb 2022)}, paper=2, qno=1, type=numerical, subtopic={Projectile motion}, answer={2.40}, answer_src=key, difficulty=2, diff_parts={K1 S1 I0 C0 F1}, diff_note={concepts: projectile speed | steps: horizontal component 9 m/s so  | vy}]
A projectile is thrown with a velocity of 18 m/s at an angle of $60^\circ$ with horizontal. The interval between the moments (in sec) when speed of the projectile is 15 m/s is ($g=10\ \mathrm{m/s^2}$)
\end{bankq}
\begin{bankq}[id=TOAS03-P2-PHY-S2-Q04, source={TOAS - Test 03 (27 Feb 2022)}, paper=2, qno=4, type=numerical, subtopic={Projectile motion}, answer={25.00}, answer_src=key, difficulty=4, diff_parts={K1 S2 I1 C0 F1}, diff_note={concepts: projectile through points with 45 degree slope | steps: roof diagonal gives horizontal separation 25 m; slopes +-1 at P and Q fix the apex; height 12.5 m and trajectory equation; range | I: symmetry about the apex}]
A commander fires a shell at certain angle of projection from `A' which clears the building (cuboid) of dimensions $20\times15\times12.5\ \mathrm{m^3}$ in its diagonal plane. Shell just clears the roof diagonally at points $P$ and $Q$ and falls on the other side of the building at $B$, and makes $45^\circ$ with horizontal at $P$ and $Q$ as shown in the figure. If the range $AB$ of the shell is $N\times\sqrt3$ metre then find $N$. ($g=10\ \mathrm{m/s^2}$)
\bfig{168.9}{TOAS03_p2_phy_s2_q04_1.png}
\end{bankq}
\begin{bankpassage}[id=TOAS07-P1-PHY-S3-G1-2, source={TOAS - Test 07 (01 May 2022)}, questions={1-2}]
A projectile is thrown from a point `O' on the ground at an angle $\theta$ from horizontal with initial velocity $v_0$ at $t=0$. At the highest point of its trajectory, the projectile splits into two parts having mass in the ratio 1 : 2. After splitting, smaller part falls vertically downwards with initial velocity zero. The other part `$t$' second after splitting falls on the ground at a distance $x$ meter from the point of projection. (Consider $v_0=20$ m/s, $\theta=37^\circ$ \& acceleration due to gravity $g=10$ m/s$^2$)
\bfig{106.1}{TOAS07_p1_phy_s3_para1_2_1.png}
\end{bankpassage}
\begin{bankq}[id=TOAS07-P1-PHY-S3-Q01, source={TOAS - Test 07 (01 May 2022)}, paper=1, qno=1, type=numerical, subtopic={Projectile motion}, answer={1.20}, answer_src=key, difficulty=2, diff_parts={K1 S1 I0 C0 F1}, diff_note={concepts: projectile problem from the passage | steps: setup of the motion; time value}, group=TOAS07-P1-PHY-S3-G1-2]
The value of $t$ is\underline{\hspace{1cm}}
\end{bankq}
\begin{bankq}[id=TOAS07-P1-PHY-S3-Q02, source={TOAS - Test 07 (01 May 2022)}, paper=1, qno=2, type=numerical, subtopic={Projectile motion}, answer={48.00}, answer_src=key, difficulty=2, diff_parts={K1 S1 I0 C0 F1}, diff_note={concepts: projectile problem from the passage | steps: setup of the motion; value of x}, group=TOAS07-P1-PHY-S3-G1-2]
The value of $x$ is \underline{\hspace{1.2cm}}
\end{bankq}
\begin{bankq}[id=TOAS19-P1-PHY-S1-Q04, source={TOAS - Test 19 (05 Feb 2023)}, paper=1, qno=4, type=multi, subtopic={1D and 2D motion}, answer={B,C,D}, answer_src=key, difficulty=3, diff_parts={K1 S2 I0 C0 F1}, diff_note={concepts: kinematics with variable acceleration | steps: velocity by integration; speed and direction; tangential acceleration; radius of curvature}]
A particle starting from rest moves with time varying acceleration $\vec a=3\hat i+4t\hat j$ m/s$^2$
\begin{qoptions}
\opt{A}{Speed of particle after 1 sec should be 5 m/s.}
\opt{B}{At $t=\dfrac32$ sec particle will be moving at $45^\circ$ with $x$ axis.}
\opt{C}{At $t=1$ sec rate of change of speed will be $\dfrac{17}{\sqrt{13}}$ m/s$^2$}
\opt{D}{At $t=1$ sec radius of curvature of the path will be $\dfrac{13}6\sqrt{13}$ m}
\end{qoptions}
\end{bankq}
\begin{bankq}[id=TOAS19-P1-PHY-S3-Q06, source={TOAS - Test 19 (05 Feb 2023)}, paper=1, qno=6, type=numerical, subtopic={River-boat problems}, answer={3.00}, answer_src=key, difficulty=7, diff_parts={K1 S3 I2 C1 F1}, diff_note={concepts: pursuit curve in a river | steps: take the water frame; equal-speed pursuit equation; integrate in polar form; drift at the bank; read m and n}]
A man crosses a river of width $d$. Flow speed of water is $v$. Speed of swimmer relative to water is also $v$. Man always heads towards the point exactly opposite to the starting point at the other bank (relative to water). If total drift is $\dfrac{md}n$. Find minimum possible value of $m+n$. ($m$ \& $n$ are integers.)
\end{bankq}

\section*{TAAS papers}
\begin{bankq}[id=TAAS01-P1-PHY-S2-Q08, source={TAAS - Test 01 (18 Jun 2023)}, paper=1, qno=8, type=multi, subtopic={Projectile motion}, answer={B}, answer_src=key, difficulty=2, diff_parts={K1 S1 I0 C0 F1}, diff_note={concepts: possible projectile paths | steps: projectile paths under uniform gravity; check the curves}]
Shown are four curves. Which amongst them can be possible trajectories traced by a projectile lunched from earth surface? (neglect air friction)
\begin{qoptions}
\opt{A}{\bfigi{78.7}{TAAS01_p1_phy_s2_q08_A.png}}
\opt{B}{\bfigi{74.7}{TAAS01_p1_phy_s2_q08_B.png}}
\opt{C}{\bfigi{56.6}{TAAS01_p1_phy_s2_q08_C.png}}
\opt{D}{\bfigi{87.4}{TAAS01_p1_phy_s2_q08_D.png}}
\end{qoptions}
\end{bankq}

\section*{TOPS papers}
\begin{bankpassage}[id=TOPS01-P1-PHY-S2-G11-12, source={TOPS - Test 01 (31 Jul 2022)}, questions={11-12}]
When two rigid bodies move over each other, their intersection point might not be moving with the velocity of either of them. For example, consider two rods $AB$ and $CD$, connected by pin $P$ that moves in slots made in the rods.
\bfig{81.8}{TOPS01_p1_phy_s2_para11_12_1.png}
If the rod $AB$ is moving with velocity $v_1$ and rod $CD$ is moving with velocity $v_2$, their intersection point (pin $P$) moves as shown:
\bfig{102.4}{TOPS01_p1_phy_s2_para11_12_2.png}
Where $A'B'$ and $C'D'$ are the new positions of the rods, while $P'$ is the new position of the pin $P$. Here the pin $P$ is having a velocity that is the sum of the velocities of the two rods. But in more complicated shapes of the sliders, this direct approach might not work, and calculus is required to calculate the velocity or acceleration of the intersection point.

In the given diagram, there is an elliptical slot as shown. The equation of the ellipse is $\dfrac{x^2}{4}+y^2=1$, where $x$ and $y$ are the coordinates in same unit. There is a slider with two pins $A$ and $B$, which are free to move along the slider so that they always remain on the elliptical slot. The slider is moved horizontally with constant speed 10 m/s. Answer the following questions.
\bfig{143.3}{TOPS01_p1_phy_s2_para11_12_3.png}
\end{bankpassage}
\begin{bankq}[id=TOPS01-P1-PHY-S2-Q11, source={TOPS - Test 01 (31 Jul 2022)}, paper=1, qno=11, type=single, subtopic={1D and 2D motion}, answer={D}, answer_src=key, difficulty=5, diff_parts={K2 S2 I1 C1 F0}, diff_note={concepts: relative velocity, ellipse geometry | steps: slope of the ellipse at the pin; velocity components along the slot; resolve with the slider speed; speed}, group=TOPS01-P1-PHY-S2-G11-12]
When $x$-coordinate of the slider is 1 m, the speed of Pin-$B$ is?
\begin{qoptions}
\opt{A}{Zero}
\opt{B}{10 m/s}
\opt{C}{2.9 m/s}
\opt{D}{10.4 m/s}
\end{qoptions}
\end{bankq}
\begin{bankq}[id=TOPS01-P1-PHY-S2-Q12, source={TOPS - Test 01 (31 Jul 2022)}, paper=1, qno=12, type=single, subtopic={1D and 2D motion}, answer={C}, answer_src=key, difficulty=9, diff_parts={K3 S3 I2 C2 F0}, diff_note={concepts: ellipse geometry, differentiation, relative acceleration | steps: position of the pin; velocity components; differentiate again; curvature acceleration; combine; magnitude | calculus on a constraint}, group=TOPS01-P1-PHY-S2-G11-12]
When $x$-coordinate of the slider is 1 m, the magnitude of acceleration of Pin-$B$ is?
\begin{qoptions}
\opt{A}{Zero}
\opt{B}{38.4 cm/s$^2$}
\opt{C}{38.5 m/s$^2$}
\opt{D}{42 cm/s$^2$}
\end{qoptions}
\end{bankq}
\begin{bankq}[id=TOPS02-P1-PHY-S1-Q01, source={TOPS - Test 02 (28 Aug 2022)}, paper=1, qno=1, type=single, subtopic={1D and 2D motion}, answer={D}, answer_src=key, difficulty=6, diff_parts={K2 S3 I1 C1 F0}, diff_note={concepts: velocity from the position vector, turning points | steps: differentiate for velocity; factor to find where each component vanishes; find the turning time; integrate speed on each leg; add the distances | vector components}]
A particle is moving such that its position vector changes with time as
\par\noindent\hfill$\begin{aligned}\vec r&=\left(36t^2-4t^3-96t\right)\hat i+\left(12t^3-108t^2+288t-17\right)\hat j\\&\quad+\left(51+3t^3-27t^2+72t\right)\hat k\end{aligned}$\hfill\null\par
What is the distance travelled by the particle in the first six seconds?
\begin{qoptions}
\opt{A}{216}
\opt{B}{384}
\opt{C}{468}
\opt{D}{572}
\end{qoptions}
\end{bankq}
\begin{bankq}[id=TOPS02-P1-PHY-S1-Q06, source={TOPS - Test 02 (28 Aug 2022)}, paper=1, qno=6, type=single, subtopic={1D and 2D motion}, answer={B}, answer_src=key, difficulty=4, diff_parts={K1 S2 I1 C1 F0}, diff_note={concepts: vector projection | steps: direction of AB; projection of AC; foot of the perpendicular; evaluate the expression}]
Consider a straight line passing through two points $A$ and $B$. Position vectors of point $A$ and $B$ are given by $\vec r_A$ and $\vec r_B$ respectively. There is another point $C$ with position vector $\vec r_C$. If a perpendicular is dropped on the line passing through $A$ and $B$ from the point $C$. Then foot of perpendicular on straight line is given by point $D$ with a position vector $\vec r_D=x_0\hat i+y_0\hat j+z_0\hat k$ if $\vec r_A=4\hat i+7\hat j+\hat k$, $\vec r_B=3\hat i+5\hat j+3\hat k$ and $\vec r_C=\hat i+3\hat k$. Then the value of $[x_0+z_0-y_0]$ will be?
\begin{qoptions}
\opt{A}{12}
\opt{B}{5}
\opt{C}{8}
\opt{D}{0}
\end{qoptions}
\end{bankq}
\begin{bankq}[id=TOPS02-P1-PHY-S2-Q10, source={TOPS - Test 02 (28 Aug 2022)}, paper=1, qno=10, type=multi, subtopic={Projectile motion}, answer={A,D}, answer_src=key, difficulty=9, diff_parts={K3 S3 I2 C1 F1}, diff_note={concepts: projectile motion, relative velocity, circular motion | steps: velocity of the train at P; velocity of the stone; time to the top; velocity at the top; position at the top | train and stone in 3D}]
A small toy train is moving with a constant speed of 10 m/s in a circle of radius $\dfrac{16}{\pi}$ m. The plane of the circle lies in horizontal $x$--$y$ plane. At time $t=0$, train is at point $P$ and moving in counter--clockwise direction. At this instant, a stone is thrown from the train with speed 10 m/s relative to train towards negative $x$--axis at an angle of $37^\circ$ with vertical $z$--axis.

(Take $g=10$ m/s$^2$, $\sin37^\circ=\dfrac35$) (Given : centre of circle at origin $0,0$)
\bfig{108.3}{TOPS02_p1_phy_s2_q10_1.png}
\begin{qoptions}
\opt{A}{Velocity of particle at highest point will become $\left(-6\hat i+10\hat j\right)$ m/s.}
\opt{B}{Velocity of particle at highest point will become $\left(-8\hat i+10\hat j\right)$ m/s.}
\opt{C}{Coordinates at highest point are, (0.3m, 7.0m, 3.2 m).}
\opt{D}{Coordinates at highest point are, (0.3m, 8.0m, 3.2m).}
\end{qoptions}
\end{bankq}
\begin{bankq}[id=TOPS02-P1-PHY-S3-Q01, source={TOPS - Test 02 (28 Aug 2022)}, paper=1, qno=1, type=numerical, subtopic={Relative motion}, answer={1.70}, answer_src=key, difficulty=7, diff_parts={K2 S2 I2 C1 F1}, diff_note={concepts: relative motion, pursuit curve | steps: rate of approach along the line of sight; component of the other velocity; angle from the figure; time to converge | chase problem}]
Point `$A$' moves uniformly with speed $v_1=20$ m/s, so that vector $\vec v_1$ is continuously `aimed' at point `$B$' which in turn moves rectilinearly and uniformly with vector $v_2=10$ m/s along the path $P\to Q$ as shown in the figure. Find the time (in sec) when the points $A$ and $B$ converge.
\bfig{137.4}{TOPS02_p1_phy_s3_q01_1.png}
\end{bankq}
\begin{bankq}[id=TOPS02-P1-PHY-S3-Q02, source={TOPS - Test 02 (28 Aug 2022)}, paper=1, qno=2, type=numerical, subtopic={Projectile motion}, answer={20.20}, answer_src=key, difficulty=10, diff_parts={K3 S3 I2 C2 F1}, diff_note={concepts: projectile on an incline, relative motion, kinematics | steps: components along and perpendicular to the incline; time to reach height 2 m; the box position at that time; the second pass; width | moving box}]
A particle is projected up an incline at $t=0$ (inclination $=53^\circ$) with speed 14 m/s and at an angle $30^\circ$ from the incline. At the same moment a box starts sliding down, the box has 2 small hole at height of 2 m on opposite walls. The particle enters form 1$^{\text{st}}$ hole and leave through the other. Find the width of the box.
\bfig{111.2}{TOPS02_p1_phy_s3_q02_1.png}
\end{bankq}
\begin{bankq}[id=TOPS02-P2-PHY-S1-Q03, source={TOPS - Test 02 (28 Aug 2022)}, paper=2, qno=3, type=multi, subtopic={Projectile motion}, answer={C}, answer_src=key, difficulty=10, diff_parts={K3 S3 I2 C2 F1}, diff_note={concepts: drag force, differential equation, projectile | steps: horizontal equation; solve the horizontal velocity; vertical equation; integrate for vy; compare the ascent and descent times | linear drag}]
A projectile of mass $m$ is fired at an angle $\theta_0$ with horizontal and initial speed $V_0$, as shown. If the air develops a drag force (resistance force) on the projectile which is proportional to its instantaneous velocity i.e. $F=-KV$ where $K$ is a positive constant of appropriate dimensions. (Assume that gravitational acceleration is constant throughout the motion i.e. $g$ m/s$^2$)
\bfig{108.3}{TOPS02_p2_phy_s1_q03_1.png}
\begin{qoptions}
\opt{A}{The time of ascent is greater than the time of decent.}
\opt{B}{Vertical ($Y$-component) component of velocity of projectile as a function of time is $\dfrac mK\left[\left(\dfrac Km V_0\sin\theta_0-g\right)e^{\frac{-K}mt}+g\right]$}
\opt{C}{Horizontal ($X$-component) component of velocity of projectile as a function of time is $V_0\cos\theta_0e^{-\frac kmt}$}
\opt{D}{Vertical ($Y$-component) component of velocity of projectile as a function of time is $\dfrac mK\left[\left(\dfrac Km V_0\sin\theta_0+g\right)e^{\frac{-2K}mt}+g\right]$}
\end{qoptions}
\end{bankq}
\begin{bankq}[id=TOPS02-P2-PHY-S1-Q06, source={TOPS - Test 02 (28 Aug 2022)}, paper=2, qno=6, type=multi, subtopic={1D and 2D motion}, answer={B,C}, answer_src=key, difficulty=7, diff_parts={K2 S2 I2 C1 F1}, diff_note={concepts: polar components, circular motion | steps: geometry of the position; radial and transverse velocity; ratio; acceleration check | velocity components about a point}]
When a particle is in motion along a curved path, its velocity at a point $A$ is along the tangent to the path at point $A$ as shown in the first figure. We may split this velocity into two components; one along position vector $OA$ and call it radial velocity and another along the perpendicular to $OA$ and call it transverse velocity. Same can be done to acceleration at $A$ and call the components radial acceleration and transverse acceleration.
\bfig{97.7}{TOPS02_p2_phy_s1_q06_1.png}
A particle is executing uniform circular motion of angular velocity $\omega$ in $x$--$y$ plane with radius $R$ center $(2R,0)$. When it is at point $A$, answer the correct option(s).
\bfig{103.2}{TOPS02_p2_phy_s1_q06_2.png}
\begin{qoptions}
\opt{A}{Ratio of angular velocity to the radial velocity of the particle about origin is $1/R$.}
\opt{B}{Ratio of angular velocity to the radial velocity of the particle about origin is $\dfrac1{2\sqrt5R}$.}
\opt{C}{Net acceleration of particle about origin is $R\omega^2$.}
\opt{D}{To find angular acceleration about the origin in the given question we can apply $a_t=\alpha r$, where $a_t$ is transverse acceleration about point $O$ and $r$ is distance between point $O$ and $A$.}
\end{qoptions}
\end{bankq}
\begin{bankq}[id=TOPS02-P2-PHY-S2-Q01, source={TOPS - Test 02 (28 Aug 2022)}, paper=2, qno=1, type=numerical, subtopic={1D and 2D motion}, answer={3.60}, answer_src=key, difficulty=6, diff_parts={K2 S2 I1 C1 F1}, diff_note={concepts: linear drag, impulse | steps: integrate the equation over the whole flight; the displacement is zero; solve for the total time}]
A ball is projected vertically upwards with a velocity 20 m/s returns back on the ground with a velocity 16 m/s. if the air resistance is proportional to the speed of ball find the air time of the ball. Acceleration of free ball is 10 SI unit.
\end{bankq}
\begin{bankq}[id=TOPS02-P2-PHY-S3-Q05, source={TOPS - Test 02 (28 Aug 2022)}, paper=2, qno=5, type=numerical, subtopic={Projectile motion}, answer={8.00}, answer_src=key, difficulty=7, diff_parts={K2 S2 I2 C1 F1}, diff_note={concepts: projectile motion, free fall | steps: the shot aimed at the monkey; time from the position of the shot; height of the monkey from the free fall; solve for n over 4 | monkey and hunter}]
A fireman fires his shot from origin aiming to monkey ($M$), location at $(12, 16)$, who falls at the time of shot. At the instant shot has passes its highest position, it is located on line $x=8$ m, the monkey is located at position $(12,n/3)$. Find the value of $n/4$. (Coordinates are in m)
\end{bankq}
\begin{bankq}[id=TOPS03-P1-PHY-S2-Q09, source={TOPS - Test 03 (20 Nov 2022)}, paper=1, qno=9, type=multi, subtopic={1D and 2D motion}, answer={A,B,C}, answer_src=key, difficulty=5, diff_parts={K1 S2 I1 C1 F1}, diff_note={concepts: related rates | steps: geometry of the ladder; differentiate the constraint; rates of area and angle | 13 12 5 triangle}]
A 13 ft ladder is leaning against the wall when its base starts to slide away. When at some instant of time the base is 12ft from the wall, then it was seen to be moving at a rate of 5 ft/sec. Select correct option(s) at the same instant of time
\bfig{178.0}{TOPS03_p1_phy_s2_q09_1.png}
\begin{qoptions}
\opt{A}{The top of the ladder is sliding down the wall at a rate of magnitude 12 ft/s.}
\opt{B}{The area of the triangle formed by the ladder, wall and ground is changing at a rate of magnitude 60 ft$^2$/s.}
\opt{C}{The angle $\theta$ between the ladder and the ground is changing at a rate of magnitude 1 rad/s.}
\opt{D}{The angle $\theta$ between the ladder and the ground is changing at a rate of magnitude 4 rad/s.}
\end{qoptions}
\end{bankq}
\begin{bankq}[id=TOPS03-P1-PHY-S2-Q10, source={TOPS - Test 03 (20 Nov 2022)}, paper=1, qno=10, type=multi, subtopic={Projectile motion}, answer={A,C,D}, answer_src=key, difficulty=5, diff_parts={K1 S2 I1 C1 F1}, diff_note={concepts: projectile motion on an incline | steps: components along and perpendicular to the plane; time of flight; ranges for both cases; compare | up and down the incline}]
Two balls are thrown from an inclined plane at angle of projection $\alpha$ with the plane one up the incline plane and other down the incline as shown in the figure. If $R_1$ \& $R_2$ be their respective ranges, then [here $T_1$ \& $T_2$ are times of flight in the two cases respectively]
\bfig{100.4}{TOPS03_p1_phy_s2_q10_1.png}
\bfig{101.5}{TOPS03_p1_phy_s2_q10_2.png}
\begin{qoptions}
\opt{A}{$h_1=h_2$}
\opt{B}{$R_2-R_1=T_1^2$}
\opt{C}{$R_2-R_1=g\sin\theta\,T_2^2$}
\opt{D}{$R_2-R_1=g\sin\theta\,T_1^2$}
\end{qoptions}
\end{bankq}
\begin{bankq}[id=TOPS04-P1-PHY-S1-Q04, source={TOPS - Test 04 (18 Dec 2022)}, paper=1, qno=4, type=multi, subtopic={1D and 2D motion}, answer={A,D}, answer_src=key, difficulty=7, diff_parts={K2 S2 I2 C1 F1}, diff_note={concepts: velocity components, curvature | steps: vy as a function of y; acceleration; path equation; force and radius of curvature}]
A particle of mass `$m$' starts from origin and move in $x$-$y$ plane. Velocity of particle is given by $\vec v=k_1\hat i+k_2\sqrt{1-\alpha y}\,\hat j$

where $k_1$, $k_2$ \& $\alpha$ are constant and `$y$' is $y$-coordinate of position of particle. Choose the correct statement(s) :
\begin{qoptions}
\opt{A}{Average force acting on particle $\dfrac{m\alpha k_2^2}2$}
\opt{B}{Force acting on particle decreases first then increases}
\opt{C}{Path of particle will be hyperbolic.}
\opt{D}{Radius of curvature of the path followed by particle will decrease first then increase.}
\end{qoptions}
\end{bankq}
\begin{bankq}[id=TOPS04-P2-PHY-S1-Q10, source={TOPS - Test 04 (18 Dec 2022)}, paper=2, qno=10, type=single, subtopic={1D and 2D motion}, answer={A}, answer_src=key, difficulty=4, diff_parts={K1 S2 I1 C1 F0}, diff_note={concepts: radius of curvature | steps: normal component of the acceleration; speed; radius}]
A particle is moving in an isolated $x$--$y$ plane. At an instant, the particle has velocity $(4\hat i+4\hat j)$ m/s and acceleration $(3\hat i+5\hat j)$ m/s$^2$. At that instant what will be the radius of curvature of its path ?
\begin{qoptions}
\opt{A}{$16\sqrt2$ m}
\opt{B}{15 m}
\opt{C}{18 m}
\opt{D}{None of these}
\end{qoptions}
\end{bankq}
\begin{bankq}[id=TOPS04-P2-PHY-S1-Q11, source={TOPS - Test 04 (18 Dec 2022)}, paper=2, qno=11, type=single, subtopic={Relative motion}, answer={C}, answer_src=key, difficulty=6, diff_parts={K2 S2 I2 C1 F0}, diff_note={concepts: relative velocity, inclines | steps: relative velocity on the first incline; on the second; solve; speed}]
There are two incline planes $AO$ and $OB$ inclined at $45^\circ$ and $60^\circ$ respectively with the horizontal as shown in the figure. When a man moves on the inclined plane $AO$, he observes that the rain drops are falling at an angle $45^\circ$ with the vertical. When the man moves on the inclined plane $OB$, he observes that the raindrops are falling vertically downward with the horizontal Then actual speed of rain is :-
\bfig{182.6}{TOPS04_p2_phy_s1_q11_1.png}
\begin{qoptions}
\opt{A}{$5\sqrt2$ m/s}
\opt{B}{$20\sqrt3$ m/s}
\opt{C}{$10\sqrt2$ m/s}
\opt{D}{None of these}
\end{qoptions}
\end{bankq}
\begin{bankq}[id=TOPS04-P2-PHY-S1-Q12, source={TOPS - Test 04 (18 Dec 2022)}, paper=2, qno=12, type=single, subtopic={Relative motion}, answer={D}, answer_src=key, difficulty=6, diff_parts={K2 S2 I2 C1 F0}, diff_note={concepts: constraint, relative velocity | steps: wedge velocities; contact conditions; cylinder velocity components; speed}]
System is shown in the figure. Assume that cylinder remain in contact with the two wedges. The speed of cylinder is :
\bfig{108.3}{TOPS04_p2_phy_s1_q12_1.png}
\begin{qoptions}
\opt{A}{$\sqrt{19-4\sqrt3}\,\dfrac u2$ m/s}
\opt{B}{$\dfrac{\sqrt{13}\,u}2$ m/s}
\opt{C}{$\sqrt3u$ m/s}
\opt{D}{$\sqrt7u$ m/s}
\end{qoptions}
\end{bankq}
\begin{bankq}[id=TOPS04-P2-PHY-S1-Q16, source={TOPS - Test 04 (18 Dec 2022)}, paper=2, qno=16, type=single, subtopic={Projectile motion}, answer={D}, answer_src=key, difficulty=6, diff_parts={K2 S2 I2 C1 F0}, diff_note={concepts: projectile, ceiling limit | steps: range and height; ceiling constraint; maximise; separation}]
Two boys ``play catch'' with a ball in a long hallway. The ceiling height is $H$, and the ball is thrown and caught at shoulder height, which we call $h$ for each boy. If the boys are capable of throwing the ball with velocity $v_0$, at what maximum separation can they play ?
\begin{qoptions}
\opt{A}{$4\sqrt{(H-h)\left(\dfrac{v_0^2}{2g}+H+h\right)}$}
\opt{B}{$4\sqrt{(H-h)\dfrac{v_0^2}{2g}+(H-h)}$}
\opt{C}{$4\sqrt{(H-h)\left(\dfrac{v_0^2}{2g}-(H+h)\right)}$}
\opt{D}{$4\sqrt{(H-h)\left(\dfrac{v_0^2}{2g}-(H-h)\right)}$}
\end{qoptions}
\end{bankq}
\begin{bankq}[id=TOPS04-P2-PHY-S2-Q08, source={TOPS - Test 04 (18 Dec 2022)}, paper=2, qno=8, type=numerical, subtopic={1D and 2D motion}, answer={5}, answer_src=solved, working={In polar form $\dot r=-5+3\cos\phi$, $r\dot\phi=-3\sin\phi$, so $r\sin\phi=16\tan^{5/3}(\phi/2)$. With $u=\tan(\phi/2)$, $t=\tfrac83\int_0^1u^{-1/3}(1+u^2)\,du=5$ s}, difficulty=9, diff_parts={K2 S3 I2 C2 F1}, diff_note={concepts: pursuit, polar coordinates | steps: radial and angular equations; eliminate time; integrate; substitute; time}]
A particle $P$ is initially at a distance $d=16$ m from a fixed point $O$. The particle $P$ moves with a velocity $\vec v=5\hat{PO}+3\hat i$. Where $\hat{PO}$ is a unit vector from $P$ to $O$ at any time $t$. Initially $\hat{PO}$ is perpendicular to $\hat i$. Find the time in seconds after which point $P$ meets point $O$.
\end{bankq}
\begin{bankq}[id=TOPS05-P1-PHY-S3-Q01, source={TOPS - Test 05 (14 Jan 2023)}, paper=1, qno=1, type=numerical, subtopic={1D and 2D motion}, answer={15.00}, answer_src=key, difficulty=5, diff_parts={K1 S2 I1 C1 F1}, diff_note={concepts: uniformly accelerated motion | steps: average velocity relation; solve for the acceleration time; uniform time}]
A car starts moving rectilinearly, first with acceleration $\alpha=5$ m/s$^2$ (the initial velocity is equal to zero), then uniformly, and finally, decelerating at the same rate $\alpha$ comes to a stop. The total time of motion equals $\tau=25$ s. The average velocity during that time is equal to $\langle v\rangle=72$ km/h. How much time (in sec) does the car move uniformly ?
\end{bankq}
\begin{bankq}[id=TOPS05-P2-PHY-S2-Q08, source={TOPS - Test 05 (14 Jan 2023)}, paper=2, qno=8, type=multi, subtopic={1D and 2D motion}, answer={A,B}, answer_src=key, difficulty=7, diff_parts={K2 S2 I2 C1 F1}, diff_note={concepts: velocity position graph, acceleration | steps: slope of the graph; acceleration as v dv dx; force; position in time}]
A body of mass 2kg is moving on a straight line for which velocity-position ($v$-$x$) graph as shown in the figure. At $t=0$, body is at origin.
\bfig{108.2}{TOPS05_p2_phy_s2_q08_1.png}
\begin{qoptions}
\opt{A}{Force acting on the body as $x=10$ m is 20 N.}
\opt{B}{$x=20\left(1-e^{-t}\right)$ m}
\opt{C}{Acceleration of the body is uniform.}
\opt{D}{Acceleration-position graph is parabola.}
\end{qoptions}
\end{bankq}

\section*{Star Review (Other) papers}
\begin{bankq}[id=STAR01-P1-PHY-S4-Q07, source={Star Review (Other) - Test 01 (12 Sep 2022)}, paper=1, qno=7, type=numerical, subtopic={1D and 2D motion}, answer={30}, answer_src=key, difficulty=5, diff_parts={K1 S2 I1 C1 F1}, diff_note={concepts: speed from an acceleration versus position graph | steps: area under the graph; change of kinetic energy per mass; final speed}]
A particle moves in rectilinear motion. The acceleration-position curve is shown below. If initial velocity of particle is 20 m/s. Velocity of particle (in m/s) after the particle has travelled 50 m is
\bfig{137.6}{STAR01_p1_phy_s4_q07_1.png}
\end{bankq}
\begin{bankq}[id=STAR01-P1-PHY-S4-Q08, source={Star Review (Other) - Test 01 (12 Sep 2022)}, paper=1, qno=8, type=numerical, subtopic={River-boat problems}, answer={20}, answer_src=key, difficulty=6, diff_parts={K1 S2 I2 C1 F1}, diff_note={concepts: drift of a swimmer in a river with varying current | steps: time to cross; position dependent drift velocity; integrate; drift}]
Velocity of river flow grows in proportion to the distance from the bank and reaches its maximum velocity of 2 m/s in the middle. Near the bank, velocity is zero. The velocity of a swimmer in still water is 5 m/s and is directed perpendicular to river flow. The width of the river is 100 m. Find the drift (in m) of the swimmer.
\end{bankq}
\begin{bankq}[id=STAR01-P2-PHY-S1-Q14, source={Star Review (Other) - Test 01 (12 Sep 2022)}, paper=2, qno=14, type=single, subtopic={Projectile motion}, answer={C}, answer_src=key, difficulty=4, diff_parts={K1 S2 I1 C1 F0}, diff_note={concepts: radius of curvature of a projectile path | steps: velocity components at the time; normal component of gravity; radius}]
A body is thrown with a velocity of 20 m/s at an angle of $60^\circ$ with horizontal. The radius of curvature of its trajectory at $t=\sqrt3$s after the body began to move is :
\begin{qoptions}
\opt{A}{0 m}
\opt{B}{5 m}
\opt{C}{10 m}
\opt{D}{15 m}
\end{qoptions}
\end{bankq}
\begin{bankq}[id=STAR01-P2-PHY-S1-Q15, source={Star Review (Other) - Test 01 (12 Sep 2022)}, paper=2, qno=15, type=single, subtopic={Projectile motion}, answer={D}, answer_src=key, difficulty=7, diff_parts={K1 S3 I2 C2 F0}, diff_note={concepts: return of a ball after leaving a circular surface | steps: condition at the top; projectile range; speed at the top; energy conservation; initial speed}]
A person rolls a small ball with speed $u$ along the floor from point A. If $x=3R$, determine the required speed $u$, so that the ball returns to A after rolling on the circular surface in the vertical plane from B to C and becoming a projectile at `C' :
\bfig{111.2}{STAR01_p2_phy_s1_q15_1.png}
\begin{qoptions}
\opt{A}{$2\sqrt{gR}$}
\opt{B}{$\sqrt{5gR}$}
\opt{C}{$\sqrt{6gR}$}
\opt{D}{$\dfrac52\sqrt{gR}$}
\end{qoptions}
\end{bankq}
\begin{bankq}[id=STAR01-P2-PHY-S1-Q20, source={Star Review (Other) - Test 01 (12 Sep 2022)}, paper=2, qno=20, type=single, subtopic={1D and 2D motion}, answer={C}, answer_src=key, difficulty=4, diff_parts={K1 S2 I1 C1 F0}, diff_note={concepts: velocity of a particle with acceleration depending on speed | steps: separate variables; integrate; form of the solution}]
The acceleration of a particle moving along $x$-axis varies with velocity $(v)$ as $a=3-2v$. If initially body was at rest, then :
\begin{qoptions}
\opt{A}{Velocity varies with time linearly.}
\opt{B}{Velocity varies with time parabolically.}
\opt{C}{Velocity varies with time exponentially.}
\opt{D}{None of these}
\end{qoptions}
\end{bankq}
\begin{bankq}[id=STAR01-P2-PHY-S2-Q01, source={Star Review (Other) - Test 01 (12 Sep 2022)}, paper=2, qno=1, type=numerical, subtopic={1D and 2D motion}, answer={4}, answer_src=key, difficulty=5, diff_parts={K1 S2 I1 C1 F1}, diff_note={concepts: acceleration from a velocity versus position graph | steps: slope of the graph; velocity at the point; acceleration equals v dv dx}]
The graph shows velocity v/s $x$-coordinate of a particle moving along $x$-axis. Find out acceleration of particle in m/s$^2$ at $x=4$m.
\bfig{144.9}{STAR01_p2_phy_s2_q01_1.png}
\end{bankq}
\begin{bankq}[id=STAR01-P2-PHY-S2-Q02, source={Star Review (Other) - Test 01 (12 Sep 2022)}, paper=2, qno=2, type=numerical, subtopic={1D and 2D motion}, answer={4}, answer_src=key, difficulty=5, diff_parts={K1 S2 I1 C1 F1}, diff_note={concepts: overtaking of a car by a bike | steps: position equations; set equal; solve the quadratic; relevant root}]
A bike moving with uniform speed of 10 m/s is at a distance of 24m from a car, the car starts from rest and moves with a uniform acceleration of 2 m/s$^2$ away from the bike. If car starts at $t=0$, the time (in sec) at which bike will overtake the car is
\end{bankq}
\begin{bankq}[id=STAR01-P2-PHY-S2-Q03, source={Star Review (Other) - Test 01 (12 Sep 2022)}, paper=2, qno=3, type=numerical, subtopic={Projectile motion}, answer={5}, answer_src=key, difficulty=5, diff_parts={K1 S2 I1 C1 F1}, diff_note={concepts: radius of curvature of a projectile | steps: velocity components; acceleration component normal; radius; scale}]
A ball is projected at $t=0$ s with velocity 20 m/s at an angle $60^\circ$ with vertical. The radius of curvature of the trajectory at $t=1$ s is $R$. Neglect air resistance and take $g=10$ m/s$^2$. The value of $\dfrac R6$ is
\end{bankq}
\begin{bankq}[id=STAR01-P2-PHY-S2-Q06, source={Star Review (Other) - Test 01 (12 Sep 2022)}, paper=2, qno=6, type=numerical, subtopic={1D and 2D motion}, answer={17}, answer_src=key, difficulty=5, diff_parts={K1 S2 I1 C1 F1}, diff_note={concepts: angular velocity of the velocity vector | steps: angle with the horizontal; derivative with time; evaluate}]
A particle is thrown horizontally with an initial speed of 5 m/s at $t=0$. If the angular velocity of its velocity vector at time $t=2$ is $\dfrac2n$, then find the value of `$n$'. [Take : $g=10$ m/s$^2$ \& neglect air drag]
\end{bankq}
\begin{bankq}[id=STAR02-P1-PHY-S3-Q01, source={Star Review (Other) - Test 02 (10 Oct 2022)}, paper=1, qno=1, type=numerical, subtopic={Projectile motion}, answer={8.99}, answer_src=key, difficulty=8, diff_parts={K1 S3 I2 C2 F1}, diff_note={concepts: ice block sliding off a roof | steps: speed at the edge; direction of the speed; time of fall; horizontal range; total distance}]
A block of ice sliding down from the top of the inclined roof of a house (angle of inclination of roof $=30^\circ$ with the horizontal). The highest and lowest points of roof are at heights of 8.1 m and 5.6 m respectively from the ground. At what horizontal distance (in m) from the starting point will the block hit the ground ? Neglect friction.
\end{bankq}
\begin{bankq}[id=STAR02-P2-PHY-S3-Q05, source={Star Review (Other) - Test 02 (10 Oct 2022)}, paper=2, qno=5, type=numerical, subtopic={Projectile motion}, answer={4}, answer_src=key, difficulty=8, diff_parts={K1 S3 I2 C2 F1}, diff_note={concepts: projectile with velocity perpendicular to the initial one | steps: components of the velocity; dot product zero; time; height; product}]
From the top of a tower of height $65$ m, a ball is projected with velocity $40$ m/s at an angle of elevation of $53^\circ$. The time (in seconds) after which its velocity becomes perpendicular to its initial velocity (take $g=10\ \mathrm{m/s^2}$) is $t$ and its height from the ground at that instant is $h$, find the value of $\dfrac{ht}{125}$.
\end{bankq}

\section*{Nurture (General) papers}
\begin{bankq}[id=NUR04-P1-PHY-S2-Q05, source={Nurture (General) - Test 04 (11 Feb 2024)}, paper=1, qno=5, type=multi, subtopic={1D and 2D motion}, answer={A,C,D}, answer_src=solved, working={Total distance 500 m gives accelerating time 5 s, speed 25 m/s, uniform stretch 15 s, accelerating distance 62.5 m, uniform distance 375 m}, difficulty=5, diff_parts={K1 S2 I1 C1 F1}, diff_note={concepts: three phase motion of a car | steps: average speed to distance; solve the accelerating time; check each statement}]
A car starts moving rectilinearly, first with an acceleration of 5 ms$^{-2}$ (initial velocity zero), then moves uniformly and finally decelerating at the same rate till it stops. The total time of journey is 25 s. The average speed during the interval is $72\,\mathrm{km\,h}^{-1}$. Which of the following statements are correct.
\begin{qoptions}
\opt{A}{Total distance travelled by the particle with uniform acceleration is 62.5 m}
\opt{B}{Total distance travelled by the particle with uniform velocity is 250 m}
\opt{C}{Time duration for which car accelerates is 5 sec}
\opt{D}{The duration for which the car moves with uniform velocity is 15 s}
\end{qoptions}
\end{bankq}
\begin{bankq}[id=NUR04-P1-PHY-S2-Q09, source={Nurture (General) - Test 04 (11 Feb 2024)}, paper=1, qno=9, type=multi, subtopic={Relative motion}, answer={C,D}, answer_src=solved, working={Rain components 2 and 2; the man needs 4 m/s to flip the horizontal component; v\^{}2=2as gives 16 m}, difficulty=5, diff_parts={K1 S2 I1 C1 F1}, diff_note={concepts: rain seen by an accelerating man | steps: rain components; condition for 45 degrees again; speed; distance}]
A man is standing on a road and observes that rain is falling at angle 45$^\circ$ with the vertical. The man starts running on the road with constant acceleration 0.5 m/s$^2$. After a certain time from the start of the motion, it appears to him that rain is still falling at angle 45$^\circ$ with the vertical, with speed $2\sqrt2$ m / s. Motion of the man is in the same vertical plane in which the rain is falling. Then which of the following statement(s) are true
\begin{qoptions}
\opt{A}{It is not possible}
\opt{B}{Speed of the rain relative to the ground is 2 m/s}
\opt{C}{Speed of the man when he finds rain to be falling at angle 45$^\circ$ with the vertical, is 4 m/s}
\opt{D}{The man has travelled a distance 16 m on the road by the time he again finds rain to be falling at angle 45$^\circ$}
\end{qoptions}
\end{bankq}
\begin{bankq}[id=NUR04-P1-PHY-S3-Q13, source={Nurture (General) - Test 04 (11 Feb 2024)}, paper=1, qno=13, type=numerical, subtopic={1D and 2D motion}, answer={4}, answer_src=solved, working={12 m at 10 km/h with a=0.5 m/s2 takes 3.33 s: rocks at 0,1,2,3 s}, difficulty=5, diff_parts={K1 S2 I1 C1 F1}, diff_note={concepts: rocks falling on an accelerating car | steps: speed conversion; distance relation; time of passage; count rocks}]
A slowly moving flat car is 12.0 m long passing a point at straight road at 10 kmph. A boy beside the road near that point tosses rocks on to the moving flat car at the rate of one per second. How many rocks will fall onto that if the train begins to accelerate at $0.5\,\mathrm{m/s^2}$ just as the first rock hits the car?
\end{bankq}
\end{document}
